\subsubsection{Subproceso 3.B: Confirmación y Registro Transaccional Oficial}

Para el levantamiento de información del Subproceso 3.B se emplearon de forma complementaria dos técnicas: la \textbf{revisión documental} y la \textbf{observación directa con registro de campo}.

La revisión documental permitió identificar el marco normativo que rige el cierre de la matrícula: el Reglamento de Matrícula (RR-0570-2022), en particular su Capítulo VIII (arts. 75° a 80°), que regula la verificación, las rectificaciones, el Cierre de Matrícula, los padrones y los reportes externos; la ficha TUPA del procedimiento de Matrícula Regular (PA5400954D, RR-3698-2024); y el marco general de la Ley Universitaria N.\textsuperscript{o} 30220. Esta técnica fue adecuada porque las reglas y plazos del subproceso se encuentran formalizados en instrumentos normativos de acceso institucional, lo que garantiza la precisión y trazabilidad de la información obtenida.

La observación directa consistió en el seguimiento del comportamiento del portal Intranet DIRCE durante el proceso, complementada con el registro de campo del periodo 2026-2 aportado por un integrante del equipo: anotaciones descriptivas del cierre del proceso que documentan la comunicación de la Facultad a la DIRCE al concluir la matrícula, la falta de acceso de la Facultad para realizar modificaciones y el desfase de las nóminas durante la primera semana. Estas técnicas fueron pertinentes porque el subproceso opera casi en su totalidad de forma mediada por la plataforma transaccional, y porque sus fricciones principales —desfases de nóminas, demoras de visado y reprocesos manuales— solo se evidencian en la operación real del periodo, no en los documentos normativos. El registro de campo permitió, además, reconstruir la cadena causal entre los problemas de los subprocesos previos (adeudos no distinguidos, fichas generadas con la base del periodo anterior, cambios de plan no convalidados) y sus efectos sobre el cierre.
